% !TEX root = ../main.tex
\chapter{Desenvolvimento}
\label{cap:desenvolvimento}

Este capítulo serve também de \emph{cábula}: mostra como fazer cada elemento.
Apagar ou substituir à vontade.

\section{Figuras}

As imagens ficam em \texttt{imagens/} e não é preciso escrever a pasta nem a
extensão. Referência automática: \cref{fig:exemplo}.

\begin{figure}[H]
  \centering
  \includegraphics[width=0.6\linewidth]{example-image-a} % ex.: {topologia}
  \caption{Legenda da figura}
  \label{fig:exemplo}
\end{figure}

Duas figuras lado a lado (\cref{fig:lado-a,fig:lado-b}):

\begin{figure}[H]
  \centering
  \begin{subfigure}{0.45\linewidth}
    \centering
    \includegraphics[width=\linewidth]{example-image-b}
    \caption{Antes}
    \label{fig:lado-a}
  \end{subfigure}
  \hfill
  \begin{subfigure}{0.45\linewidth}
    \centering
    \includegraphics[width=\linewidth]{example-image-c}
    \caption{Depois}
    \label{fig:lado-b}
  \end{subfigure}
  \caption{Comparação antes/depois}
\end{figure}

\section{Tabelas}

\begin{table}[H]
  \centering
  \caption{Endereçamento IP}   % nas tabelas a legenda vai em cima
  \label{tab:enderecos}
  \begin{tabular}{llll}
    \toprule
    \textbf{Router} & \textbf{Interface} & \textbf{IP privado} & \textbf{IP público} \\
    \midrule
    R1 & Lo0  & 10.1.1.1/32 & --- \\
       & g1/0 & 10.1.2.1/30 & --- \\
    R2 & Lo1  & ---         & 47.73.239.22/32 \\
    \bottomrule
  \end{tabular}
\end{table}

Para tabelas que ocupam várias páginas usar \texttt{longtable}.

\section{Código}

\begin{lstlisting}[language=Java, caption={Exemplo de Web Service}, label={lst:exemplo}]
@WebService
public class Servico {
    @WebMethod
    public String ola(String nome) {
        return "Olá, " + nome + "! Já funcionam acentos.";
    }
}
\end{lstlisting}

Código num ficheiro externo:
\verb|\lstinputlisting[language=Python]{codigo/script.py}|.
Configuração de um router, sem linguagem: \verb|\begin{lstlisting}...\end{lstlisting}|.

\section{Equações}

\begin{equation}
  \label{eq:exemplo}
  T = \frac{L}{R} = \frac{\SI{1500}{\byte}}{\SI{100}{\mega\bit\per\second}}
\end{equation}

Ver \cref{eq:exemplo}. Os acrónimos já usados aparecem curtos: \gls{bgp}, \gls{as}.
Livro de referência:~\cite{kurose2021}.
